\documentclass[10pt]{article}
\usepackage[utf8]{inputenc}
\usepackage[T1]{fontenc}
\usepackage{lmodern}
\usepackage[paperwidth=384pt,paperheight=900pt,margin=10pt]{geometry}
\usepackage{amsmath,amssymb}
\usepackage{array}
\usepackage{enumitem}
\usepackage{xcolor}
\usepackage{microtype}

\setlength{\parindent}{0pt}
\setlength{\parskip}{4pt}
\setlist[itemize]{leftmargin=12pt}
\setlist[enumerate]{leftmargin=14pt}

\sloppy
\allowdisplaybreaks
\emergencystretch=2em

\begin{document}

\section*{Uppgift 2}

\subsection*{Original question}

En maskin producerar skruvar vars längd måste
understiga 2 cm för att vara användbara i
produktionen. Från erfarenhet vet man att skruvarnas
längd kan anses vara normalfördelad med väntevärde
\(\mu=1.8\) och standardavvikelse \(\sigma=0.4\).

\begin{enumerate}[label=\alph*)]
\item Vad är sannolikheten att en skruv är användbar
i produktionen?

\item Skruvarna förpackas i lådor om 10 stycken.
Vad är sannolikheten att alla skruvar i en låda är
användbara?

\item Hur många skruvar måste produceras för att
sannolikheten att minst en av de producerade
skruvarna är användbar är större än 0.99?

\item På sistone har man börjat misstänka att
maskinen producerar med ett väntevärde större än
1.8. Däremot har man ingen anledning att ifrågasätta
den antagna standardavvikelsen.

I ett experiment med 10 oberoende mätningar av
skruvlängder har man fått att skruvlängdernas
medelvärde är 2.0.

Testa hypotesen att maskinens väntevärde är 1.8 i
ett ensidigt test (d.v.s att man bara förkastar
hypotesen om det finns tydlig anledning att tro att
väntevärdet är större än 1.8). Genomför testet på
5\% signifikansnivå. \((8p)\)
\end{enumerate}

\subsection*{Compiled notes from Yacoub + Obid}
% UNCERTAIN: The supplied note images are not visibly labelled by author,
% so Yacoub's and Obid's notes are compiled together as the provided notes.
% UNCERTAIN: One handwritten setup looks like sigma = 0.5, but the
% original question says sigma = 0.4 and the later hypothesis-test notes
% also use sigma = 0.4. The calculation z = 0.5 matches sigma = 0.4.

\subsubsection*{What we are doing}

The screw length is normally distributed.
A screw is usable if its length is less than 2 cm.

First we find the probability that one screw is
usable. Then we use powers for repeated independent
screws. Finally we do a one-sided \(Z\)-test because
\(\sigma\) is known.

\subsubsection*{Given information}

\[
X\sim N(\mu,\sigma)
=N(1.8,0.4).
\]

\[
\mu=1.8,
\qquad
\sigma=0.4.
\]

Usable screw:
\[
X<2.
\]

\subsubsection*{Step-by-step solution}

\textbf{a) Probability that one screw is usable}

The notes write:
\[
P(X<2).
\]

First standardize:
\[
Z=\frac{X-\mu}{\sigma}.
\]

Substitute:
\begin{align*}
z
&=
\frac{2-1.8}{0.4}
\\[2mm]
&=
\frac{0.2}{0.4}
\\[2mm]
&=0.5.
\end{align*}

So:
\begin{align*}
P(X<2)
&=
P(Z<0.5)
\\[2mm]
&=
\Phi(0.5).
\end{align*}

From the \(Z\)-table:
\[
\Phi(0.5)=0.6915.
\]

Therefore:
\[
P(X<2)=0.6915.
\]

The notes write this as about \(69\%\).

\textbf{b) All 10 screws are usable}

En maskin producerar skruvar vars längd måste
understiga 2 cm för att vara användbara i
produktionen. Från erfarenhet vet man att skruvarnas
längd kan anses vara normalfördelad med väntevärde
\(\mu=1.8\) och standardavvikelse \(\sigma=0.4\).

\begin{enumerate}[label=\alph*)]
\item Vad är sannolikheten att en skruv är användbar
i produktionen?

\item Skruvarna förpackas i lådor om 10 stycken.
Vad är sannolikheten att alla skruvar i en låda är
användbara?

\item Hur många skruvar måste produceras för att
sannolikheten att minst en av de producerade
skruvarna är användbar är större än 0.99?

\item På sistone har man börjat misstänka att
maskinen producerar med ett väntevärde större än
1.8. Däremot har man ingen anledning att ifrågasätta
den antagna standardavvikelsen.

I ett experiment med 10 oberoende mätningar av
skruvlängder har man fått att skruvlängdernas
medelvärde är 2.0.

Testa hypotesen att maskinens väntevärde är 1.8 i
ett ensidigt test (d.v.s att man bara förkastar
hypotesen om det finns tydlig anledning att tro att
väntevärdet är större än 1.8). Genomför testet på
5\% signifikansnivå. \((8p)\)
\end{enumerate}

From part a), the chance that one screw is usable is:
\[
p=0.6915.
\]

There are 10 screws in one box.
The notes use:
\[
P(\text{10 usable})
=p^{10}.
\]

So:
\begin{align*}
P(\text{10 usable})
&=
(0.6915)^{10}
\\[2mm]
&\approx
0.0248.
\end{align*}

\textbf{c) At least one usable screw}

En maskin producerar skruvar vars längd måste
understiga 2 cm för att vara användbara i
produktionen. Från erfarenhet vet man att skruvarnas
längd kan anses vara normalfördelad med väntevärde
\(\mu=1.8\) och standardavvikelse \(\sigma=0.4\).

\begin{enumerate}[label=\alph*)]
\item Vad är sannolikheten att en skruv är användbar
i produktionen?

\item Skruvarna förpackas i lådor om 10 stycken.
Vad är sannolikheten att alla skruvar i en låda är
användbara?

\item Hur många skruvar måste produceras för att
sannolikheten att minst en av de producerade
skruvarna är användbar är större än 0.99?

\item På sistone har man börjat misstänka att
maskinen producerar med ett väntevärde större än
1.8. Däremot har man ingen anledning att ifrågasätta
den antagna standardavvikelsen.

I ett experiment med 10 oberoende mätningar av
skruvlängder har man fått att skruvlängdernas
medelvärde är 2.0.

Testa hypotesen att maskinens väntevärde är 1.8 i
ett ensidigt test (d.v.s att man bara förkastar
hypotesen om det finns tydlig anledning att tro att
väntevärdet är större än 1.8). Genomför testet på
5\% signifikansnivå. \((8p)\)
\end{enumerate}

We want:
\[
P(\text{at least one usable})>0.99.
\]

The notes say it is easier to use the opposite event:
\[
P(\text{at least one usable})
=
1-P(\text{all are not usable}).
\]

The chance that one screw is bad is:
\begin{align*}
1-0.6915
&=
0.3085.
\end{align*}

So the chance that \(n\) screws are all bad is:
\[
(0.3085)^n.
\]

Thus:
\[
1-(0.3085)^n>0.99.
\]

Solve for \(n\):
\begin{align*}
1-(0.3085)^n
&>0.99
\\[2mm]
-(0.3085)^n
&>-0.01
\\[2mm]
(0.3085)^n
&<0.01.
\end{align*}

Take logarithms:
\begin{align*}
\ln\left((0.3085)^n\right)
&<
\ln(0.01)
\\[2mm]
n\ln(0.3085)
&<
\ln(0.01).
\end{align*}

Since \(\ln(0.3085)<0\), the inequality reverses
when dividing:
\begin{align*}
n
&>
\frac{\ln(0.01)}{\ln(0.3085)}
\\[2mm]
&\approx
3.915.
\end{align*}

The notes round up because half a screw does not make
sense:
\[
n=4.
\]

\textbf{d) One-sided hypothesis test}

En maskin producerar skruvar vars längd måste
understiga 2 cm för att vara användbara i
produktionen. Från erfarenhet vet man att skruvarnas
längd kan anses vara normalfördelad med väntevärde
\(\mu=1.8\) och standardavvikelse \(\sigma=0.4\).

\begin{enumerate}[label=\alph*)]
\item Vad är sannolikheten att en skruv är användbar
i produktionen?

\item Skruvarna förpackas i lådor om 10 stycken.
Vad är sannolikheten att alla skruvar i en låda är
användbara?

\item Hur många skruvar måste produceras för att
sannolikheten att minst en av de producerade
skruvarna är användbar är större än 0.99?

\item På sistone har man börjat misstänka att
maskinen producerar med ett väntevärde större än
1.8. Däremot har man ingen anledning att ifrågasätta
den antagna standardavvikelsen.

I ett experiment med 10 oberoende mätningar av
skruvlängder har man fått att skruvlängdernas
medelvärde är 2.0.

Testa hypotesen att maskinens väntevärde är 1.8 i
ett ensidigt test (d.v.s att man bara förkastar
hypotesen om det finns tydlig anledning att tro att
väntevärdet är större än 1.8). Genomför testet på
5\% signifikansnivå. \((8p)\)
\end{enumerate}

The notes list:
\[
n=10,
\qquad
\bar{x}=2.0,
\qquad
\sigma=0.4.
\]

The significance level is:
\[
\alpha=0.05.
\]

Set up the hypotheses:
\[
H_0:\mu=1.8.
\]

\[
H_1:\mu>1.8.
\]

This is a one-sided right-tailed test.

The notes use the critical value:
\[
z_{0.05}=1.645.
\]

Now calculate the test statistic:
\[
Z=
\frac{\bar{x}-\mu_0}{\sigma/\sqrt{n}}.
\]

Substitute:
\begin{align*}
Z
&=
\frac{2.0-1.8}{0.4/\sqrt{10}}
\\[2mm]
&\approx
1.581.
\end{align*}

Compare with the critical value:
\[
1.581<1.645.
\]

So the notes conclude that \(H_0\) is still correct.

\textbf{English clarification:}
In hypothesis-test language, this means we do not
reject \(H_0\) at the 5\% significance level.
There is not enough evidence that the mean is greater
than \(1.8\).

\subsubsection*{Final answer}

\begin{enumerate}[label=\alph*)]
\item
\[
P(X<2)=0.6915.
\]

\item
\[
(0.6915)^{10}\approx 0.0248.
\]

\item
\[
n=4.
\]

\item
\[
Z\approx 1.581.
\]

Since:
\[
1.581<1.645,
\]
we do not reject \(H_0:\mu=1.8\).
\end{enumerate}

\subsubsection*{What to remember}

For normal probabilities:
\[
Z=\frac{X-\mu}{\sigma}.
\]

For all independent screws to be usable:
\[
p^{10}.
\]

For at least one:
\[
1-P(\text{none}).
\]

For a one-sided \(Z\)-test with known \(\sigma\):
\[
Z=
\frac{\bar{x}-\mu_0}{\sigma/\sqrt{n}}.
\]

\section*{Uppgift 3}

\subsection*{Original question}

Betrakta funktionen

\[
f(x)=
\begin{cases}
ax^3,
& 0<x<1,\\
0,
& \text{annars}.
\end{cases}
\]

\begin{enumerate}[label=\alph*)]
\item Bestäm konstanten \(a\) sådan att \(f\) är en
täthetsfunktion för en stokastisk variabel \(X\).

\item Bestäm fördelningsfunktionen för \(X\).

\item Beräkna
\[
P(|X|<\tfrac12).
\]

\item Anta att \(X\) bara antar värden \(x\) med
\[
0<x<1.
\]
Bestäm täthetsfunktionen för
\[
Y=\ln X.
\]
\((6p)\)
\end{enumerate}

\subsection*{Compiled notes from Yacoub + Obid}
% UNCERTAIN: The supplied note images are not visibly labelled by author,
% so Yacoub's and Obid's notes are compiled together as the provided notes.

\subsubsection*{What we are doing}

First we choose \(a\) so that the total area under
the density is \(1\).

Then we find the CDF \(F(x)\). After that, we use the
CDF to calculate a probability. Finally we transform
\[
Y=\ln X
\]
by first finding the CDF of \(Y\), then differentiating.

\subsubsection*{Given information}

\[
f(x)=
\begin{cases}
ax^3,
& 0<x<1,\\
0,
& \text{annars}.
\end{cases}
\]

For a density:
\[
\int_{-\infty}^{\infty}f(x)\,dx=1.
\]

\subsubsection*{Step-by-step solution}

\textbf{a) Find the constant \(a\)}

The notes say that a PDF must integrate to \(1\).
Since the density is only nonzero for \(0<x<1\):
\[
\int_0^1 ax^3\,dx=1.
\]

Integrate:
\begin{align*}
\int_0^1 ax^3\,dx
&=
\left[
a\frac{x^4}{4}
\right]_0^1
\\[2mm]
&=
a\left(
\frac{1^4}{4}
-
\frac{0^4}{4}
\right)
\\[2mm]
&=
\frac{a}{4}.
\end{align*}

Set equal to \(1\):
\begin{align*}
\frac{a}{4}
&=1
\\[2mm]
a&=4.
\end{align*}

So:
\[
a=4.
\]

\textbf{b) Find the CDF}

Betrakta funktionen

\[
f(x)=
\begin{cases}
ax^3,
& 0<x<1,\\
0,
& \text{annars}.
\end{cases}
\]

\begin{enumerate}[label=\alph*)]
\item Bestäm konstanten \(a\) sådan att \(f\) är en
täthetsfunktion för en stokastisk variabel \(X\).

\item Bestäm fördelningsfunktionen för \(X\).

\item Beräkna
\[
P(|X|<\tfrac12).
\]

\item Anta att \(X\) bara antar värden \(x\) med
\[
0<x<1.
\]
Bestäm täthetsfunktionen för
\[
Y=\ln X.
\]
\((6p)\)
\end{enumerate}

Now we know:
\[
f(x)=4x^3,
\qquad 0<x<1.
\]

For \(0<x<1\), the notes integrate from \(0\) to
\(x\):
\begin{align*}
F(x)
&=
\int_0^x 4t^3\,dt
\\[2mm]
&=
\left[
t^4
\right]_0^x
\\[2mm]
&=
x^4-0
\\[2mm]
&=
x^4.
\end{align*}

Therefore the full CDF is:
\[
F(x)=
\begin{cases}
0,
& x\le 0,\\
x^4,
& 0<x<1,\\
1,
& x\ge 1.
\end{cases}
\]

\textbf{English clarification:}
The notes explain this as: before the density starts,
the probability is \(0\). After the whole support has
passed, the probability is \(1\).

\textbf{c) Calculate \(P(|X|<\frac12)\)}

Betrakta funktionen

\[
f(x)=
\begin{cases}
ax^3,
& 0<x<1,\\
0,
& \text{annars}.
\end{cases}
\]

\begin{enumerate}[label=\alph*)]
\item Bestäm konstanten \(a\) sådan att \(f\) är en
täthetsfunktion för en stokastisk variabel \(X\).

\item Bestäm fördelningsfunktionen för \(X\).

\item Beräkna
\[
P(|X|<\tfrac12).
\]

\item Anta att \(X\) bara antar värden \(x\) med
\[
0<x<1.
\]
Bestäm täthetsfunktionen för
\[
Y=\ln X.
\]
\((6p)\)
\end{enumerate}

The notes rewrite:
\[
|X|<\frac12
\]
as:
\[
-\frac12<X<\frac12.
\]

So:
\[
P(|X|<\tfrac12)
=
P\left(
-\frac12<X<\frac12
\right).
\]

Using the CDF:
\[
P\left(
-\frac12<X<\frac12
\right)
=
F\left(\frac12\right)
-
F\left(-\frac12\right).
\]

From part b):
\[
F\left(\frac12\right)
=
\left(\frac12\right)^4.
\]

And since \(-\frac12\le 0\):
\[
F\left(-\frac12\right)=0.
\]

Therefore:
\begin{align*}
P(|X|<\tfrac12)
&=
\left(\frac12\right)^4-0
\\[2mm]
&=
\frac{1}{16}
\\[2mm]
&=
0.0625.
\end{align*}

The notes write this as:
\[
6.25\%.
\]

\textbf{d) Density for \(Y=\ln X\)}

Betrakta funktionen

\[
f(x)=
\begin{cases}
ax^3,
& 0<x<1,\\
0,
& \text{annars}.
\end{cases}
\]

\begin{enumerate}[label=\alph*)]
\item Bestäm konstanten \(a\) sådan att \(f\) är en
täthetsfunktion för en stokastisk variabel \(X\).

\item Bestäm fördelningsfunktionen för \(X\).

\item Beräkna
\[
P(|X|<\tfrac12).
\]

\item Anta att \(X\) bara antar värden \(x\) med
\[
0<x<1.
\]
Bestäm täthetsfunktionen för
\[
Y=\ln X.
\]
\((6p)\)
\end{enumerate}

The notes start with the CDF:
\[
F_Y(y)=P(Y\le y).
\]

Since:
\[
Y=\ln X,
\]
we get:
\[
F_Y(y)=P(\ln X\le y).
\]

Put both sides as powers of \(e\):
\[
\ln X\le y
\quad \Longleftrightarrow \quad
X\le e^y.
\]

So:
\[
F_Y(y)=P(X\le e^y).
\]

Use the CDF from part b):
\[
P(X\le e^y)=F_X(e^y).
\]

For \(e^y\) to fall inside \(0<x<1\), we need:
\[
y<0.
\]

For \(y<0\):
\begin{align*}
F_Y(y)
&=
F_X(e^y)
\\[2mm]
&=
(e^y)^4
\\[2mm]
&=
e^{4y}.
\end{align*}

Thus the CDF of \(Y\) is:
\[
F_Y(y)=
\begin{cases}
e^{4y},
& y<0,\\
1,
& y\ge 0.
\end{cases}
\]

\textbf{English clarification:}
As \(y\to -\infty\), \(e^{4y}\to 0\), so this also
matches the lower tail.

Now differentiate the CDF to get the density:
\begin{align*}
f_Y(y)
&=
\frac{d}{dy}F_Y(y)
\\[2mm]
&=
\frac{d}{dy}e^{4y}
\\[2mm]
&=
4e^{4y}.
\end{align*}

So:
\[
f_Y(y)=
\begin{cases}
4e^{4y},
& y<0,\\
0,
& \text{annars}.
\end{cases}
\]

\subsubsection*{Final answer}

\begin{enumerate}[label=\alph*)]
\item
\[
a=4.
\]

\item
\[
F(x)=
\begin{cases}
0,
& x\le 0,\\
x^4,
& 0<x<1,\\
1,
& x\ge 1.
\end{cases}
\]

\item
\[
P(|X|<\tfrac12)
=
\frac{1}{16}
=0.0625.
\]

\item
\[
f_Y(y)=
\begin{cases}
4e^{4y},
& y<0,\\
0,
& \text{annars}.
\end{cases}
\]
\end{enumerate}

\subsubsection*{What to remember}

For a density:
\[
\int f(x)\,dx=1.
\]

For a CDF:
\[
F(x)=P(X\le x).
\]

For intervals:
\[
P(a<X<b)=F(b)-F(a).
\]

For \(Y=\ln X\), use the CDF method:
\[
F_Y(y)=P(Y\le y).
\]

\section*{Uppgift 4}

\subsection*{Original question}

De stokastiska variabler \(X,Y\) är oberoende.
\(X\) är likformigt fördelad på intervallet
\((-1,1)\) och \(Y\):s fördelning är symmetrisk
kring origo.

\begin{enumerate}[label=\alph*)]
\item Beräkna
\[
P(-0.1<X<0.2).
\]

\item Bestäm
\[
E(2X-1)
\quad\text{och}\quad
V(2X-1).
\]

\item Bestäm fördelningsfunktionen för \(2X-1\).

\item Beräkna
\[
P(\max(X,Y)\le 0).
\]
\((6p)\)
\end{enumerate}

\subsection*{Compiled notes from Yacoub + Obid}
% UNCERTAIN: The supplied note images are not visibly labelled by author,
% so Yacoub's and Obid's notes are compiled together as the provided notes.

\subsubsection*{What we are doing}

We use that \(X\) is uniform on \((-1,1)\).
That gives a constant density and simple formulas
for mean and variance.

We also use that \(Y\) is symmetric around \(0\), so:
\[
P(Y\le 0)=0.5.
\]

Finally, since \(X\) and \(Y\) are independent, we can
multiply probabilities in part d).

\subsubsection*{Given information}

\[
X\sim U(-1,1).
\]

\[
X \text{ and } Y
\text{ are independent.}
\]

\[
Y \text{ is symmetric around }0.
\]

For \(X\sim U(a,b)\), the notes use:
\[
f_X(x)=\frac{1}{b-a}.
\]

\[
E(X)=\frac{a+b}{2}.
\]

\[
V(X)=\frac{(b-a)^2}{12}.
\]

\subsubsection*{Step-by-step solution}

\textbf{Collect facts about \(X\)}

For \(X\sim U(-1,1)\):
\[
f_X(x)
=
\frac{1}{1-(-1)}
=
\frac12,
\qquad -1<x<1.
\]

Mean:
\begin{align*}
E(X)
&=
\frac{-1+1}{2}
\\[2mm]
&=0.
\end{align*}

Variance:
\begin{align*}
V(X)
&=
\frac{(1-(-1))^2}{12}
\\[2mm]
&=
\frac{4}{12}
\\[2mm]
&=
\frac13.
\end{align*}

Because \(X\) is symmetric around \(0\):
\[
P(X\le 0)=0.5.
\]

Because \(Y\) is symmetric around \(0\):
\[
P(Y\le 0)=0.5.
\]

\textbf{a) Calculate \(P(-0.1<X<0.2)\)}

For a uniform distribution on \((-1,1)\), the notes
use:
\[
P(c<X<d)
=
\int_c^d \frac12\,dx
=
\frac{d-c}{2}.
\]

Here:
\[
c=-0.1,
\qquad
d=0.2.
\]

So:
\begin{align*}
P(-0.1<X<0.2)
&=
\frac12
\bigl(0.2-(-0.1)\bigr)
\\[2mm]
&=
\frac12(0.3)
\\[2mm]
&=
0.15.
\end{align*}

\textbf{b) Find \(E(2X-1)\)}

De stokastiska variabler \(X,Y\) är oberoende.
\(X\) är likformigt fördelad på intervallet
\((-1,1)\) och \(Y\):s fördelning är symmetrisk
kring origo.

\begin{enumerate}[label=\alph*)]
\item Beräkna
\[
P(-0.1<X<0.2).
\]

\item Bestäm
\[
E(2X-1)
\quad\text{och}\quad
V(2X-1).
\]

\item Bestäm fördelningsfunktionen för \(2X-1\).

\item Beräkna
\[
P(\max(X,Y)\le 0).
\]
\((6p)\)
\end{enumerate}

The notes use:
\[
E(aX+b)=aE(X)+b.
\]

So:
\begin{align*}
E(2X-1)
&=
2E(X)-1
\\[2mm]
&=
2\cdot 0-1
\\[2mm]
&=
-1.
\end{align*}

\textbf{b) Find \(V(2X-1)\)}

De stokastiska variabler \(X,Y\) är oberoende.
\(X\) är likformigt fördelad på intervallet
\((-1,1)\) och \(Y\):s fördelning är symmetrisk
kring origo.

\begin{enumerate}[label=\alph*)]
\item Beräkna
\[
P(-0.1<X<0.2).
\]

\item Bestäm
\[
E(2X-1)
\quad\text{och}\quad
V(2X-1).
\]

\item Bestäm fördelningsfunktionen för \(2X-1\).

\item Beräkna
\[
P(\max(X,Y)\le 0).
\]
\((6p)\)
\end{enumerate}

The notes use:
\[
V(aX+b)=a^2V(X).
\]

So:
\begin{align*}
V(2X-1)
&=
2^2V(X)
\\[2mm]
&=
4\cdot \frac13
\\[2mm]
&=
\frac43.
\end{align*}

\textbf{c) CDF for \(2X-1\)}

De stokastiska variabler \(X,Y\) är oberoende.
\(X\) är likformigt fördelad på intervallet
\((-1,1)\) och \(Y\):s fördelning är symmetrisk
kring origo.

\begin{enumerate}[label=\alph*)]
\item Beräkna
\[
P(-0.1<X<0.2).
\]

\item Bestäm
\[
E(2X-1)
\quad\text{och}\quad
V(2X-1).
\]

\item Bestäm fördelningsfunktionen för \(2X-1\).

\item Beräkna
\[
P(\max(X,Y)\le 0).
\]
\((6p)\)
\end{enumerate}

Let:
\[
Z=2X-1.
\]

The notes isolate \(X\):
\begin{align*}
F_Z(z)
&=
P(Z\le z)
\\[2mm]
&=
P(2X-1\le z)
\\[2mm]
&=
P(2X\le z+1)
\\[2mm]
&=
P\left(X\le \frac{z+1}{2}\right).
\end{align*}

Since \(X\) is between \(-1\) and \(1\), the new
variable \(Z=2X-1\) is between:
\[
2(-1)-1=-3
\]
and:
\[
2(1)-1=1.
\]

So \(Z\) is uniform on \((-3,1)\).

The notes use the uniform CDF formula:
\[
F_Z(z)
=
\frac{z-\text{lower}}{\text{upper}-\text{lower}}.
\]

Here lower \(=-3\), upper \(=1\), so:
\[
F_Z(z)
=
\frac{z-(-3)}{1-(-3)}
=
\frac{z+3}{4}.
\]

Therefore:
\[
F_{2X-1}(z)=
\begin{cases}
0,
& z\le -3,\\
\dfrac{z+3}{4},
& -3<z<1,\\
1,
& z\ge 1.
\end{cases}
\]

\textbf{d) Calculate \(P(\max(X,Y)\le 0)\)}

De stokastiska variabler \(X,Y\) är oberoende.
\(X\) är likformigt fördelad på intervallet
\((-1,1)\) och \(Y\):s fördelning är symmetrisk
kring origo.

\begin{enumerate}[label=\alph*)]
\item Beräkna
\[
P(-0.1<X<0.2).
\]

\item Bestäm
\[
E(2X-1)
\quad\text{och}\quad
V(2X-1).
\]

\item Bestäm fördelningsfunktionen för \(2X-1\).

\item Beräkna
\[
P(\max(X,Y)\le 0).
\]
\((6p)\)
\end{enumerate}

The notes rewrite:
\[
\max(X,Y)\le 0
\]
as:
\[
X\le 0
\quad\text{and}\quad
Y\le 0.
\]

So:
\[
P(\max(X,Y)\le 0)
=
P(X\le 0,\ Y\le 0).
\]

Because \(X\) and \(Y\) are independent:
\[
P(X\le 0,\ Y\le 0)
=
P(X\le 0)P(Y\le 0).
\]

Use the symmetry facts:
\[
P(X\le 0)=0.5,
\qquad
P(Y\le 0)=0.5.
\]

Therefore:
\begin{align*}
P(\max(X,Y)\le 0)
&=
0.5\cdot 0.5
\\[2mm]
&=
0.25.
\end{align*}

\subsubsection*{Final answer}

\begin{enumerate}[label=\alph*)]
\item
\[
P(-0.1<X<0.2)=0.15.
\]

\item
\[
E(2X-1)=-1.
\]

\[
V(2X-1)=\frac43.
\]

\item
\[
F_{2X-1}(z)=
\begin{cases}
0,
& z\le -3,\\
\dfrac{z+3}{4},
& -3<z<1,\\
1,
& z\ge 1.
\end{cases}
\]

\item
\[
P(\max(X,Y)\le 0)=0.25.
\]
\end{enumerate}

\subsubsection*{What to remember}

For \(X\sim U(a,b)\):
\[
f_X(x)=\frac{1}{b-a}.
\]

\[
E(X)=\frac{a+b}{2}.
\]

\[
V(X)=\frac{(b-a)^2}{12}.
\]

For a linear transformation:
\[
E(aX+b)=aE(X)+b.
\]

\[
V(aX+b)=a^2V(X).
\]

For independent events, multiply probabilities.

\section*{Uppgift 5}

\subsection*{Original question}

Den stokastiska variabeln \(X\) har täthetsfunktionen

\[
f_X(x)=
\begin{cases}
\theta(1+x)^{-1-\theta},
& x\ge 0,\\
0,
& \text{annars}.
\end{cases}
\]

Det är känt att parametern \(\theta\) är antingen
2, 3 eller 4.

Man har ett slumpmässigt stickprov från denna
fördelning omfattande två värden, nämligen
\[
0.2
\quad\text{och}\quad
0.8.
\]
Ange ML-skattningen av \(\theta\). \((6p)\)

\subsection*{Compiled notes from Yacoub + Obid}
% UNCERTAIN: The supplied note images are not visibly labelled by author,
% so Yacoub's and Obid's notes are compiled together as the provided notes.

\subsubsection*{What we are doing}

We find the likelihood for the two observations:
\[
x_1=0.2,
\qquad
x_2=0.8.
\]

The question says that \(\theta\) can only be:
\[
2,\ 3,\ 4.
\]

So the notes do not solve by differentiation.
They calculate \(L(2)\), \(L(3)\), and \(L(4)\),
then choose the largest one.

\subsubsection*{Given information}

\[
f_X(x)=
\begin{cases}
\theta(1+x)^{-1-\theta},
& x\ge 0,\\
0,
& \text{annars}.
\end{cases}
\]

\[
x_1=0.2,
\qquad
x_2=0.8.
\]

\[
\theta\in\{2,3,4\}.
\]

\subsubsection*{Step-by-step solution}

\textbf{1) Write the likelihood}

For independent observations, the notes use:
\[
L(\theta)
=
f_X(x_1)\cdot f_X(x_2).
\]

Substitute the density:
\begin{align*}
L(\theta)
&=
\left[
\theta(1+x_1)^{-1-\theta}
\right]
\left[
\theta(1+x_2)^{-1-\theta}
\right].
\end{align*}

Now use:
\[
x_1=0.2,
\qquad
x_2=0.8.
\]

Then:
\[
1+x_1=1.2,
\qquad
1+x_2=1.8.
\]

So:
\begin{align*}
L(\theta)
&=
\left[
\theta(1.2)^{-1-\theta}
\right]
\left[
\theta(1.8)^{-1-\theta}
\right]
\\[2mm]
&=
\theta^2
\left(1.2\cdot 1.8\right)^{-1-\theta}.
\end{align*}

Since:
\[
1.2\cdot 1.8=2.16,
\]
we get:
\[
L(\theta)
=
\theta^2(2.16)^{-(1+\theta)}.
\]

\textbf{2) Evaluate the possible values}

The notes calculate only the possible values
\(\theta=2,3,4\).

For \(\theta=2\):
\begin{align*}
L(2)
&=
2^2(2.16)^{-(1+2)}
\\[2mm]
&=
4(2.16)^{-3}
\\[2mm]
&\approx 0.3969.
\end{align*}

For \(\theta=3\):
\begin{align*}
L(3)
&=
3^2(2.16)^{-(1+3)}
\\[2mm]
&=
9(2.16)^{-4}
\\[2mm]
&\approx 0.4135.
\end{align*}

For \(\theta=4\):
\begin{align*}
L(4)
&=
4^2(2.16)^{-(1+4)}
\\[2mm]
&=
16(2.16)^{-5}
\\[2mm]
&\approx 0.3403.
\end{align*}

\textbf{3) Choose the largest likelihood}

Compare:
\[
L(2)\approx 0.3969,
\]
\[
L(3)\approx 0.4135,
\]
\[
L(4)\approx 0.3403.
\]

The largest is:
\[
L(3).
\]

Therefore the ML estimate is:
\[
\hat{\theta}_{ML}=3.
\]

\subsubsection*{Final answer}

\[
\hat{\theta}_{ML}=3.
\]

\subsubsection*{What to remember}

When the parameter can only take a few possible
values, evaluate the likelihood at each possible
value.

Then choose the value with the largest likelihood:
\[
\hat{\theta}_{ML}
=
\arg\max_{\theta} L(\theta).
\]

\section*{Uppgift 6}

\subsection*{Original question}

En tentamen i matematik på Mittuniversitetet skrevs
i två grupper, på grund av att alla inte hade
möjligheten att skriva tentan på campus. Grupp 1
skrev på distans, grupp 2 på campus. De möjliga
betygen var \(F=\) underkänd, \(G=\) godkänd och
\(VG=\) väl godkänd.

Tentamensresultatet var:

\[
\begin{array}{c|ccc}
& F & G & VG\\
\hline
\text{grupp 1} & 15 & 10 & 5\\
\text{grupp 2} & 23 & 27 & 10
\end{array}
\]

Testa hypotesen att betygsfördelningen är samma i de
bägge grupperna på 5\% signifikansnivå. \((4p)\)

\subsection*{Compiled notes from Yacoub + Obid}
% UNCERTAIN: The supplied note images are not visibly labelled by author,
% so Yacoub's and Obid's notes are compiled together as the provided notes.

\subsubsection*{What we are doing}

We test whether the grade distribution is the same
in both groups.

The notes use a chi-square test with observed and
expected frequencies.

\subsubsection*{Given information}

Observed frequencies:
\[
\begin{array}{c|ccc}
& F & G & VG\\
\hline
\text{group 1} & 15 & 10 & 5\\
\text{group 2} & 23 & 27 & 10
\end{array}
\]

Group totals:
\[
\text{group 1}: 15+10+5=30.
\]

\[
\text{group 2}: 23+27+10=60.
\]

Grade totals:
\[
F:15+23=38.
\]

\[
G:10+27=37.
\]

\[
VG:5+10=15.
\]

Total number of students:
\[
30+60=90.
\]

\subsubsection*{Step-by-step solution}

\textbf{1) State the hypotheses}

The notes write:
\[
H_0:
\text{both groups have the same grade distribution.}
\]

\[
H_1:
\text{they do not have the same grade distribution.}
\]

\textbf{2) Calculate expected frequencies}

If \(H_0\) is true, the expected frequency in each
cell is:
\[
E=
\frac{
(\text{group total})(\text{grade total})
}{
\text{student total}
}.
\]

For group 1:
\begin{align*}
E_{1,F}
&=
\frac{30\cdot 38}{90}
\\[2mm]
&=12.67.
\end{align*}

\begin{align*}
E_{1,G}
&=
\frac{30\cdot 37}{90}
\\[2mm]
&=12.33.
\end{align*}

\begin{align*}
E_{1,VG}
&=
\frac{30\cdot 15}{90}
\\[2mm]
&=5.00.
\end{align*}

For group 2:
\begin{align*}
E_{2,F}
&=
\frac{60\cdot 38}{90}
\\[2mm]
&=25.33.
\end{align*}

\begin{align*}
E_{2,G}
&=
\frac{60\cdot 37}{90}
\\[2mm]
&=24.67.
\end{align*}

\begin{align*}
E_{2,VG}
&=
\frac{60\cdot 15}{90}
\\[2mm]
&=10.00.
\end{align*}

\textbf{3) Calculate the chi-square statistic}

The notes use:
\[
\chi^2
=
\sum
\frac{(O-E)^2}{E}.
\]

For group 1:
\begin{align*}
\frac{(15-12.67)^2}{12.67}
&\approx 0.430.
\end{align*}

\begin{align*}
\frac{(10-12.33)^2}{12.33}
&\approx 0.440.
\end{align*}

\begin{align*}
\frac{(5-5.00)^2}{5.00}
&=0.
\end{align*}

For group 2:
\begin{align*}
\frac{(23-25.33)^2}{25.33}
&\approx 0.214.
\end{align*}

\begin{align*}
\frac{(27-24.67)^2}{24.67}
&\approx 0.220.
\end{align*}

\begin{align*}
\frac{(10-10.00)^2}{10.00}
&=0.
\end{align*}

Add all terms:
\begin{align*}
\chi^2
&\approx
0.430+0.440+0
\\
&\quad
+0.214+0.220+0
\\[2mm]
&=
1.304.
\end{align*}

\textbf{4) Degrees of freedom}

The notes use:
\[
\text{df}
=
(\text{rows}-1)(\text{columns}-1).
\]

Here:
\[
\text{rows}=2,
\qquad
\text{columns}=3.
\]

So:
\begin{align*}
\text{df}
&=
(2-1)(3-1)
\\[2mm]
&=
1\cdot 2
\\[2mm]
&=2.
\end{align*}

\textbf{5) Critical value}

At 5\% significance level:
\[
\alpha=0.05.
\]

With:
\[
\text{df}=2,
\]
the notes use the chi-square table value:
\[
\chi^2_{0.05,2}=5.991.
\]

Compare:
\[
1.304<5.991.
\]

The test statistic is less than the rejection value.
Therefore we do not reject \(H_0\).

The notes conclude that \(H_0\) is correct:
the grade distribution is equal across both groups.

\textbf{English clarification:}
In hypothesis-test language, this means we do not
have enough evidence at the 5\% level to say that
the grade distributions are different.

\subsubsection*{Final answer}

\[
\chi^2_{\text{obs}}\approx 1.304.
\]

\[
\text{df}=2.
\]

\[
\chi^2_{0.05,2}=5.991.
\]

Since:
\[
1.304<5.991,
\]
we do not reject:
\[
H_0:
\text{same grade distribution in both groups.}
\]

\subsubsection*{What to remember}

For expected frequencies in a contingency table:
\[
E=
\frac{
(\text{row total})(\text{column total})
}{
\text{grand total}
}.
\]

For a chi-square test:
\[
\chi^2
=
\sum
\frac{(O-E)^2}{E}.
\]

Degrees of freedom:
\[
\text{df}
=
(\text{rows}-1)(\text{columns}-1).
\]

If:
\[
\chi^2_{\text{obs}}
<
\chi^2_{\text{critical}},
\]
we do not reject \(H_0\).

\end{document}
